% GENERATED FILE. Edit canonical lessons, then run npm run generate:books.
% canonical-source-hash: fnv1a64:5ffd39342c96bbf6
% canonical-lessons: KA-C323-aparadha, KA-C323-bandhisu, KA-C323-thane, KA-C323-tanikhe, KA-C323-sankisu

\chapter{A Crime in the Town}
\label{ch:ka-a-crime-in-the-town}
\hlchaptermodality{323}

\begin{chapteropening}
\textbf{By the end of this chapter:} \emph{I can say a crime, to arrest, a police station, an investigation and to suspect.}

Complete the last lesson of chapter 323: I can say a crime, to arrest, a police station, an investigation and to suspect.
\end{chapteropening}

\section[\texorpdfstring{\kn{ಅಪರಾಧ} (aparādha)}{ಅಪರಾಧ (aparādha)}]{\kn{ಅಪರಾಧ} (aparādha) — a crime; an offence}

\label{lesson:KA-C323-aparadha}



\begin{quote}
Before the new one: say the Kannada for a throne, then the Kannada for a crown.
\end{quote}

\subsection*{You'll want to know: \kn{ಅಪರಾಧ}}
\textbf{\kn{ಅಪರಾಧ}} — \emph{aparādha} — \textquotedblleft{}a crime; an offence\textquotedblright{}.

\begin{quote}
\kn{ಅಪರಾಧ} \kn{ಮಾಡಿದರೆ} \kn{ಶಿಕ್ಷೆ} \kn{ಸಿಗುತ್ತದೆ}.
\end{quote}

\begin{quote}
\emph{aparādha māḍidare śikṣe siguttade.}
\end{quote}

\begin{quote}
If you commit a crime, you get punished.
\end{quote}

One who commits a crime is an \kn{ಅಪರಾಧಿ}, a criminal. A crime is \textquotedblleft{}done\textquotedblright{}: \kn{ಅಪರಾಧ} \kn{ಮಾಡು}.

\begin{cousinweb}
\textbf{\kn{ಅಪರಾಧ}} (\emph{aparādha}) is borrowed from Sanskrit \emph{aparādha}, \textquotedblleft{}an offence, a fault\textquotedblright{} (CDIAL 439).
\end{cousinweb}

\begin{grammarlens}[title={What you've built}]
One more word for telling what happened.
\end{grammarlens}

\subsection*{Your turn}
\begin{itemize}
\raggedright
  \item \emph{Say it:} \emph{aparādha}
  \item \emph{Say it:} \emph{aparādha}, once more
  \item \emph{Say it:} \emph{aparādha māḍidare śikṣe siguttade.} — If you commit a crime, you get punished.
  \item \emph{Recall:} say the Kannada for a throne, then the Kannada for a crown, then say \emph{aparādha} again
\end{itemize}

\begin{tcolorbox}[breakable,colback=teal!4,colframe=teal!35!black,title={Before you move on}]
What does \kn{ಅಪರಾಧ} mean? (\textquotedblleft{}A crime; an offence\textquotedblright{}.) Say it once more.
\end{tcolorbox}
\section[\texorpdfstring{\kn{ಬಂಧಿಸು} (bandhisu)}{ಬಂಧಿಸು (bandhisu)}]{\kn{ಬಂಧಿಸು} (bandhisu) — to arrest}

\label{lesson:KA-C323-bandhisu}



\begin{quote}
Before the new one: say the Kannada for a crown, then the Kannada for a crime.
\end{quote}

\subsection*{You'll want to know: \kn{ಬಂಧಿಸು}}
\textbf{\kn{ಬಂಧಿಸು}} — \emph{bandhisu} — \textquotedblleft{}to arrest\textquotedblright{}.

\begin{quote}
\kn{ಪೊಲೀಸರು} \kn{ಕಳ್ಳನನ್ನು} \kn{ಬೆಳಗ್ಗೆ} \kn{ಬಂಧಿಸಿದರು}.
\end{quote}

\begin{quote}
\emph{polīsaru kaḷḷanannu beḷagge bandhisidaru.}
\end{quote}

\begin{quote}
The police arrested the thief in the morning.
\end{quote}

Past (I …): \kn{ಬಂಧಿಸಿದೆ}, \emph{bandhiside}.

The person arrested takes -\kn{ಅನ್ನು}: \kn{ಕಳ್ಳನನ್ನು} \kn{ಬಂಧಿಸು}.

\begin{cousinweb}
\textbf{\kn{ಬಂಧಿಸು}} (\emph{bandhisu}) is built on Sanskrit \emph{bandha}, \textquotedblleft{}a bond\textquotedblright{}, and also \textquotedblleft{}custody\textquotedblright{} (CDIAL 9136), from the root \emph{bandh}, \textquotedblleft{}to bind\textquotedblright{}, with the Kannada verb ending -\emph{isu}.
\end{cousinweb}

\begin{grammarlens}[title={What you've built}]
One more word for telling what happened.
\end{grammarlens}

\subsection*{Your turn}
\begin{itemize}
\raggedright
  \item \emph{Say it:} \emph{bandhisu}
  \item \emph{Say it:} \emph{bandhisu}, once more
  \item \emph{Say it:} \emph{polīsaru kaḷḷanannu beḷagge bandhisidaru.} — The police arrested the thief in the morning.
  \item \emph{Recall:} say the Kannada for a crown, then the Kannada for a crime, then say \emph{bandhisu} again
\end{itemize}

\begin{tcolorbox}[breakable,colback=teal!4,colframe=teal!35!black,title={Before you move on}]
What does \kn{ಬಂಧಿಸು} mean? (\textquotedblleft{}To arrest\textquotedblright{}.) Say it once more.
\end{tcolorbox}
\section[\texorpdfstring{\kn{ಠಾಣೆ} (ṭhāṇe)}{ಠಾಣೆ (ṭhāṇe)}]{\kn{ಠಾಣೆ} (ṭhāṇe) — a police station}

\label{lesson:KA-C323-thane}



\begin{quote}
Before the new one: say the Kannada for a crime, then the Kannada for to arrest.
\end{quote}

\subsection*{You'll want to know: \kn{ಠಾಣೆ}}
\textbf{\kn{ಠಾಣೆ}} — \emph{ṭhāṇe} — \textquotedblleft{}a police station\textquotedblright{}.

\begin{quote}
\kn{ನಾನು} \kn{ಠಾಣೆಗೆ} \kn{ಹೋಗಿ} \kn{ದೂರು} \kn{ಕೊಟ್ಟೆ}.
\end{quote}

\begin{quote}
\emph{nānu ṭhāṇege hōgi dūru koṭṭe.}
\end{quote}

\begin{quote}
I went to the police station and made a complaint.
\end{quote}

In full it is \kn{ಪೊಲೀಸ್} \kn{ಠಾಣೆ}.

\begin{cousinweb}
\textbf{\kn{ಠಾಣೆ}} (\emph{ṭhāṇe}) came through Hindustani \emph{thānā}, \textquotedblleft{}a post, a police station\textquotedblright{}, which goes back to Sanskrit \emph{sthāna}, \textquotedblleft{}a place, an abode\textquotedblright{} (CDIAL 13753).
\end{cousinweb}

\begin{grammarlens}[title={What you've built}]
One more word for telling what happened.
\end{grammarlens}

\subsection*{Your turn}
\begin{itemize}
\raggedright
  \item \emph{Say it:} \emph{ṭhāṇe}
  \item \emph{Say it:} \emph{ṭhāṇe}, once more
  \item \emph{Say it:} \emph{nānu ṭhāṇege hōgi dūru koṭṭe.} — I went to the police station and made a complaint.
  \item \emph{Recall:} say the Kannada for a crime, then the Kannada for to arrest, then say \emph{ṭhāṇe} again
\end{itemize}

\begin{tcolorbox}[breakable,colback=teal!4,colframe=teal!35!black,title={Before you move on}]
What does \kn{ಠಾಣೆ} mean? (\textquotedblleft{}A police station\textquotedblright{}.) Say it once more.
\end{tcolorbox}
\section[\texorpdfstring{\kn{ತನಿಖೆ} (tanikhe)}{ತನಿಖೆ (tanikhe)}]{\kn{ತನಿಖೆ} (tanikhe) — an investigation}

\label{lesson:KA-C323-tanikhe}



\begin{quote}
Before the new one: say the Kannada for to arrest, then the Kannada for a police station.
\end{quote}

\subsection*{You'll want to know: \kn{ತನಿಖೆ}}
\textbf{\kn{ತನಿಖೆ}} — \emph{tanikhe} — \textquotedblleft{}an investigation\textquotedblright{}.

\begin{quote}
\kn{ಪೊಲೀಸರು} \kn{ಆ} \kn{ಘಟನೆಯ} \kn{ತನಿಖೆ} \kn{ಶುರು} \kn{ಮಾಡಿದರು}.
\end{quote}

\begin{quote}
\emph{polīsaru ā ghaṭaneya tanikhe śuru māḍidaru.}
\end{quote}

\begin{quote}
The police began an investigation into that incident.
\end{quote}

\kn{ತನಿಖೆ} \kn{ಮಾಡು}, to investigate. What is investigated takes the possessive -\kn{ಅ}: \kn{ಘಟನೆಯ} \kn{ತನಿಖೆ}.

\begin{cousinweb}
\textbf{\kn{ತನಿಖೆ}} (\emph{tanikhe}) is usually traced to Hindustani \emph{tanqīh}, \textquotedblleft{}an inquiry, an investigation\textquotedblright{}, a word of Arabic origin; the path into Kannada is not settled.
\end{cousinweb}

\begin{grammarlens}[title={What you've built}]
One more word for telling what happened.
\end{grammarlens}

\subsection*{Your turn}
\begin{itemize}
\raggedright
  \item \emph{Say it:} \emph{tanikhe}
  \item \emph{Say it:} \emph{tanikhe}, once more
  \item \emph{Say it:} \emph{polīsaru ā ghaṭaneya tanikhe śuru māḍidaru.} — The police began an investigation into that incident.
  \item \emph{Recall:} say the Kannada for to arrest, then the Kannada for a police station, then say \emph{tanikhe} again
\end{itemize}

\begin{tcolorbox}[breakable,colback=teal!4,colframe=teal!35!black,title={Before you move on}]
What does \kn{ತನಿಖೆ} mean? (\textquotedblleft{}An investigation\textquotedblright{}.) Say it once more.
\end{tcolorbox}
\section[\texorpdfstring{\kn{ಶಂಕಿಸು} (śaṅkisu)}{ಶಂಕಿಸು (śaṅkisu)}]{\kn{ಶಂಕಿಸು} (śaṅkisu) — to suspect}

\label{lesson:KA-C323-sankisu}



\begin{quote}
Before the new one: say the Kannada for a police station, then the Kannada for an investigation.
\end{quote}

\subsection*{You'll want to know: \kn{ಶಂಕಿಸು}}
\textbf{\kn{ಶಂಕಿಸು}} — \emph{śaṅkisu} — \textquotedblleft{}to suspect\textquotedblright{}.

\begin{quote}
\kn{ಪೊಲೀಸರು} \kn{ಆ} \kn{ಕೆಲಸಗಾರನನ್ನು} \kn{ಶಂಕಿಸಿದರು}.
\end{quote}

\begin{quote}
\emph{polīsaru ā kelasagāranannu śaṅkisidaru.}
\end{quote}

\begin{quote}
The police suspected that worker.
\end{quote}

Past (I …): \kn{ಶಂಕಿಸಿದೆ}, \emph{śaṅkiside}.

The person suspected takes -\kn{ಅನ್ನು}. The noun is \kn{ಶಂಕೆ}, a suspicion.

\begin{cousinweb}
\textbf{\kn{ಶಂಕಿಸು}} (\emph{śaṅkisu}) is built on Sanskrit \emph{śaṅkā}, \textquotedblleft{}fear, distrust, suspicion\textquotedblright{} (CDIAL 12258), with the Kannada verb ending -\emph{isu}.
\end{cousinweb}

\begin{grammarlens}[title={What you've built}]
One more word for telling what happened.
\end{grammarlens}

\subsection*{Your turn}
\begin{itemize}
\raggedright
  \item \emph{Say it:} \emph{śaṅkisu}
  \item \emph{Say it:} \emph{śaṅkisu}, once more
  \item \emph{Say it:} \emph{polīsaru ā kelasagāranannu śaṅkisidaru.} — The police suspected that worker.
  \item \emph{Recall:} say the Kannada for a police station, then the Kannada for an investigation, then say \emph{śaṅkisu} again
\end{itemize}

\begin{tcolorbox}[breakable,colback=teal!4,colframe=teal!35!black,title={Before you move on}]
What does \kn{ಶಂಕಿಸು} mean? (\textquotedblleft{}To suspect\textquotedblright{}.) Say it once more.
\end{tcolorbox}
